\subsection{An explicit cubic class that stops at intrinsic class two}

\begin{statement}{E.38}{Proved}{A finite class-three obstruction after a full class-two lift}
Let $N=\pi_1(U_{3,\overline{\mathbb Q}})^{(3)}$ and
$V=N^{\rm ab}$, with the cubic, cusp labels and arithmetic
origin of E.37. Let $\alpha_P\in H^1(\mathbb Q,V)$ be the
abelian class of the actual rational tangential section at
$P=A_0=[0:-1:1]$, with tangent coefficient two for the
parameter $s=-x/z$. Write $e_Q\in V$ for positive inertia
at $Q=B_0=[-1:0:1]$. This is a $G_{\mathbb Q}$-equivariant
copy of $\mathbb Z_3(1)$ in $V$. The class
\[
 \alpha=\alpha_P+27\kappa_3(19)e_Q
 \quad\text{has toric coordinates}\quad (2,19^{27}).
 \tag{E38.1}
\]
Its fiber in the full intrinsic class-two section map
(E37.1) is nonempty, and is the torsor (E37.12).
There is no section above $\alpha$ in any of the following
three actual arithmetic quotients:
\[
 N/\gamma_4N,\qquad N/N'',\qquad
 N/\gamma_3I_{\rm marked}.
 \tag{E38.2}
\]
Here the marked cusp subgroup is the closed normal subgroup
generated by $Q$ alone, or by any Galois-stable subset of the
cusps above the base puncture $1$ that contains $Q$.
The obstruction is detected at the rational prime $19$ by
a finite $G_{\mathbb Q}$-equivariant quotient of $N$ of
order $3^{30}$ and nilpotency class three. Thus every
class-two section above (E38.1) fails to lift through each
of (E38.2).

\emph{Manuscript: \S51.16, folios 733v--736v, scans 172--178,
(371)--(379); Part E(c),(d) and the fourth exactness follow-up.
The quotient maps below specify the precise transfer from the
full intrinsic cubic group.}
\end{statement}
\begin{proof}
\emph{The positive full class-two calculation.}
The tangential section defining $\alpha_P$ is a section of
the full arithmetic fundamental group, so it has a lift at
every quotient in question. Relative to the origin at $O$
its path class is the framed Kummer class of E.37, with
the tangent coefficient two added at $P$. Its coordinates
are unramified away from $2$ and $3$: all the unscaled
framed values (E37.7) are units away from $3$.

We apply the full $44$-dimensional criterion (E37.11),
over $K=\mathbb Q(\zeta)$. At the prime over $2$ the Kummer
class of $19$ is zero. Indeed $19\in1+2\mathcal O_{K,2}$,
and taking the third power is an automorphism of this
pro-two unit group. The coherent third-power roots give
the zero local pro-three Kummer class. Consequently
$\alpha$ is locally cohomologous to $\alpha_P$ there and
has zero local obstruction.

At either prime over $19$ the class $\alpha_P$ is unramified.
Choose tame inertia so that $\kappa_3(19)(\tau)=1$.
The inertia value in (E37.9) is then $27e_Q$.
Writing $a_P=\alpha_P(F)$ and $u=\kappa_3(19)(F)$ gives
the exact local numerator
\[
 D=(a_P+27u e_Q)\wedge27e_Q-27k_Q\wedge e_Q
   =27(a_P-k_Q)\wedge e_Q\in27M_0.
 \tag{E38.3}
\]
For these two primes the CM scalars are
$-2+3\zeta$ and $-5-3\zeta$.
The exact integral matrices (E37.6), evaluated by
\texttt{e\_intrinsic\_19.py}, have the same Smith
valuations
\[
 0^{\times8},\quad1^{\times6},\quad
 2^{\times18},\quad3^{\times12}.
 \tag{E38.4}
\]
The sum is $78$, also obtained from the determinant
formula (E37.15), and the largest is three.
Thus $27M_0\subseteq d_{\mathfrak q}(M/L)$ at each prime
over $19$. Equation (E38.3) has an integral correction.
At every other prime away from $3$ there is no inertia
in (E38.1), so its support representative is zero.
The corrections therefore have finite support and
satisfy the convergence condition of E.37. This proves
the full intrinsic class-two lift and its stated fiber.

\emph{A continuous arithmetic quotient with its exact map.}
Use the rational unit tangent $s=-x/z$ at $P$ for this
calculation. The covering coordinate is $t=s^3$.
The other covering root is
$y/z=-(1-s^3)^{1/3}$, with rational constant term $-1$.
Taking coefficientwise roots of these power series
identifies this tangential fiber with the fiber over
the unit tangent at $t=0$. It gives a continuous
$G_{\mathbb Q}$-equivariant inclusion
\[
 N\longrightarrow F=\langle x,y\rangle^{\mathrm{pro}\text{-}3},
 \qquad N=\ker(F\longrightarrow(\mathbb Z/3)^2).
 \tag{E38.5}
\]
Here $x,y$ are the positive base inertia loops at $0,1$.
A geometric path from this tangent to the tangent at
$O$ identifies the two geometric groups. Changing that
path conjugates the identification. On $V$ this
conjugation is trivial, and the change of arithmetic
origin is the actual tangential cocycle already used
to define $\alpha_P$. Thus (E38.5) acts on precisely
the class (E38.1), with its origin accounted for.

Put $I_y=\ker(F\to\mathbb Z_3)$, where $x\mapsto1$ and
$y\mapsto0$, and let $T=1+U$. The continuous quotient
by the closed commutator subgroup of $I_y$ is
\[
 R=F/[I_y,I_y]=\mathbb Z_3[[U]]\rtimes\mathbb Z_3,
 \quad (a,v)(b,w)=(a+b,v+T^a w),
 \quad x=(1,0),\ y=(0,1).
 \tag{E38.6}
\]
For an explicit inverse construction, $I_y$ is the
closed normal subgroup generated by $y$. Its
abelianization is therefore generated, as a continuous
$\mathbb Z_3[[\mathbb Z_3]]$-module, by the image of
$y$, with $T$ acting by conjugation by $x$. This gives
a continuous surjection from the completed group ring
onto $I_y^{\rm ab}$. In the other direction the free
pro-three property sends $x,y$ to $(1,0),(0,1)$ in
the displayed semidirect product. Its restriction to
$I_y$ is abelian and sends the first module generator
back to one. The composite on the completed group
ring is the identity, so the first surjection is an
isomorphism. The section generated by $x$ now gives
exactly (E38.6), including its quotient map and topology.
Its restriction to $N$ is onto
\[
 H=\{(a,v)\in R:a\in3\mathbb Z_3,\ v(1)\in3\mathbb Z_3\}
   =3\mathbb Z_3\ltimes J,
 \qquad J=(3,U)\mathbb Z_3[[U]] .
 \tag{E38.7}
\]
Indeed $F\to R$ is onto and induces the identity on
$F^{\rm ab}=\mathbb Z_3^2$, with second coordinate the
augmentation $v(1)$, so the inverse image of $H$ is
exactly $N$.

The normalized genuine Galois action used in E.20 is
$x\mapsto x^{\chi(g)}$,
$y\mapsto f_g^{-1}y^{\chi(g)}f_g$, with $f_g\in F'$.
Its image in $R$ has $f_g\in I_y/[I_y,I_y]$, which
commutes with $y$. It follows on the actual generators
and then by continuity that
\[
 g(a,v(T))=(\chi(g)a,\ \chi(g)v(T^{\chi(g)})).
 \tag{E38.8}
\]
This proves the arithmetic action and preservation of
$H$ without discarding an unknown conjugating term.

The group $R$ is metabelian, so $N''$ maps to one.
Every inertia loop above $1$ maps to a conjugate of
a power of $y$, hence into the abelian normal module
in (E38.6). The subgroup $I_{\rm marked}$ maps into
that same module, and $\gamma_3I_{\rm marked}$ maps
to one. Thus there are the actual continuous arithmetic
surjections
\[
 N/N''\longrightarrow H,\qquad
 N/\gamma_3I_{\rm marked}\longrightarrow H.
 \tag{E38.9}
\]
More precisely, put $J_y=[I_y,I_y]$, taking the closed
commutator subgroup. Since $J_y\subset F'\subset N$,
the kernel of $N\to H$ is exactly $J_y$. The kernels
of the two maps in (E38.9) are therefore
$J_y/N''$ and $J_y/\gamma_3I_{\rm marked}$,
respectively. Indeed the whole subgroup
$[I_{\rm marked},I_{\rm marked}]$ already maps to one.

\emph{The local equation, including normalization of the section.}
At $19$ fix a compatible root-of-unity identification
with $\zeta_9$ reducing to $4\in\mathbb F_{19}$.
Since $4$ has order nine and $2^{(19-1)/9}=4$, the
unramified Kummer value is
$\kappa_3(2)(F)\equiv1\pmod9$.
In (E38.5) the section at $P$ with tangent coefficient
two has first coordinate $3\kappa_3(2)$; changing the
tangent sends $t=s^3$ to a leading coefficient $2^3$.
The added cusp class has first coordinate zero and
second augmentation $3\cdot27\kappa_3(19)$, because
inertia at $Q$ maps to a conjugate of $y^3$.

For every possible section above $\alpha$, its
Frobenius first coordinate can consequently be set
equal to $a=3$ by geometric conjugation by $x^{3c}$:
the change is $3(1-19)c$, and
$(\kappa_3(2)(F)-1)/18\in\mathbb Z_3$.
This conjugation preserves the inertia augmentation.
The tame inertia value has first coordinate zero and
is $(0,v(T))$ with
\[
 v(1)=81.
 \tag{E38.10}
\]
The cyclotomic action is $19$ at Frobenius and one
on tame inertia. In (E38.6)--(E38.8), the relation
$F\tau F^{-1}=\tau^{19}$ now becomes
\[
 T^3v(T^{19})=v(T).
 \tag{E38.11}
\]
The module coordinate of Frobenius cancels because
the normal module is abelian. This equation therefore
tests every possible Frobenius lift.

For a quick exact obstruction over $\mathbb Q_3$, set
$z=\log T$. Equation (E38.11) becomes
$e^{3z}h(19z)=h(z)$, $h(0)=81$.
Multiplying by $e^{z/6}$ shows that every positive
coefficient of the resulting series vanishes, since
$19^n-1\ne0$. Thus its unique formal solution is
\[
 v(T)=81T^{-1/6}.
 \tag{E38.12}
\]
For $r=-1/6$, every $r-j$, $j\in\mathbb Z$, has
three-adic valuation $-1$. Hence
\[
 v_3\!\left(81\binom{-1/6}{n}\right)
     =4-n-v_3(n!).
\]
At $n=4$ this equals $-1$, excluding an element of
$\mathbb Z_3[[U]]$. The finite certificate below also
detects the failure after an intrinsic class-three
quotient, and makes this nonexistence independently
checkable by integer matrix multiplication.

\emph{A finite class-three certificate.}
Reduce the module to
$A=(\mathbb Z/729)[U]/(U^5)$ and the first coordinate
modulo $2187$. The binomial coefficients of
$(1+U)^{2187}-1$ through degree four are divisible by
$729$, so these choices define a finite quotient
$R_6$ of (E38.6). Formula (E38.8) descends to it:
substitution and its inverse are defined over $A$,
and changing $\chi$ by $2187$ changes neither the
substitution nor the scalar action.
Its subgroup $H_6$ is the image of (E38.7), of order
$3^{35}$.

Write $v(U)=81+\sum_{i=1}^4v_iU^i$. The four
positive-degree coefficients of (E38.11) are
$Kv=b\pmod{729}$, with
\[
 K=\begin{pmatrix}
 18&0&0&0\\
 228&360&0&0\\
 1539&7581&6858&0\\
 7315&86640&205770&130320
 \end{pmatrix},\qquad
 b=\begin{pmatrix}-243\\-243\\-81\\0\end{pmatrix}.
 \tag{E38.13}
\]
The single integral row
\[
 \ell=(1,660,558,162)
 \quad\text{satisfies}\quad
 \ell K=0\pmod{729},\qquad \ell b=486\pmod{729}.
 \tag{E38.14}
\]
It excludes all four-tuples $v_i$.

To obtain the stronger class-three assertion, put
$\delta=T^3-1$. Direct commutators in (E38.7) give
$\gamma_{c+1}H_6=\delta^cJ$ for $c\ge1$:
the module is abelian, and conjugation by the
generator $x^3$ acts by $T^3$.
In the present ring the subgroup $\gamma_4H_6$
is generated by
\[
 3\delta^3=81U^3+243U^4,\qquad
 U\delta^3=27U^4.
 \tag{E38.15}
\]
These elements generate a subgroup of order $3^5$,
so $H_6/\gamma_4H_6$ has order $3^{30}$.
Its class is exactly three: $\delta^2U=9U^3+18U^4$
is in $\gamma_3H_6$ and has $U^3$ coefficient nine,
whereas every element of (E38.15) has that coefficient
divisible by $81$.
The row $\ell$ also annihilates their coefficient
columns $(0,0,81,243)^{\rm t}$ and
$(0,0,0,27)^{\rm t}$ modulo $729$.
Consequently (E38.14) still excludes the local tame
equation after quotienting by $\gamma_4H_6$.
The composite
\[
 N\longrightarrow H\longrightarrow H_6
              \longrightarrow H_6/\gamma_4H_6
 \tag{E38.16}
\]
is continuous, onto and $G_{\mathbb Q}$-equivariant.
Functoriality of the lower central series makes it
factor through $N/\gamma_4N$. Its factorizations
through the other two groups in (E38.2) follow from
(E38.9). Any section of any one of those arithmetic
quotients above $\alpha$ would therefore give the
excluded local section. This proves all the claimed
nonlifting assertions.

The program \texttt{e\_next\_depth.py} constructs the
matrix from the power series, checks (E38.14) and
(E38.15), and tests membership in the full integer
column lattice $[K\mid3^mI]$ by Hermite normal form.
The first failed modulus is $3^6$.
Adding the columns of $\delta^cJ$ tests the same
finite quotient at each lower central stage: the
class-one and class-two systems have solutions,
whereas the class-three system has the explicit
annihilating row (E38.14). Its saved output and
JSON certificate include every matrix and residue
used here.
\end{proof}
